% ============================================================
% Breadcrumb navigation + stub status  (deterministic chrome)
% Palette (breadcrumb / breadcrumbborder / breadcrumbtitle) is defined in boxes.tex.
% Emitted by the breadcrumb generator; never hand-rolled, never LLM-drawn.
% Requires: tcolorbox, etoolbox (\ifblank), amsmath (\text); amber color.
% \input this from the preamble after boxes.tex.
% ============================================================
\tcbset{
  breadcrumbstyle/.style={
    colback=breadcrumb, colframe=breadcrumbborder, coltitle=breadcrumbtitle,
    fonttitle=\bfseries\small, arc=2pt, boxrule=0.6pt,
    left=6pt, right=6pt, top=4pt, bottom=4pt,
    before skip=6pt, after skip=8pt,
  }
}

% \breadcrumb{<chapter subject>}{<prior display>}{<current display>}{<next display>}
% Empty prior/next are omitted together with their arrows. ASCII-only content.
\NewDocumentCommand{\breadcrumb}{m m m m}{%
  \begin{tcolorbox}[breadcrumbstyle, title={#1}]%
    \centering$%
      \ifblank{#2}{}{\text{#2}\;\to\;}%
      \textbf{#3}%
      \ifblank{#4}{}{\;\to\;\text{#4}}%
    $%
  \end{tcolorbox}%
}

% \stubstatus  -- appended after the breadcrumb on planned (stub) chapters.
\NewDocumentCommand{\stubstatus}{}{%
  \begin{tcolorbox}[breadcrumbstyle, colback=amber!10, colframe=amber]%
    \textbf{Status:} Planned%
  \end{tcolorbox}%
}

% ============================================================
% Chapter roadmap exposition (replaces the verbose "Structural Roadmap").
% Simple, templated: where the development moves from -> to, and what it studies.
% \chapterroadmap{<prior topic>}{<next topic>}{<topics studied, prose or list>}
% prior/next are the registry neighbours (same source as \breadcrumb); the topic
% list is the chapter's section inventory -- so this can be emitted deterministically.
% ============================================================
\NewDocumentCommand{\chapterroadmap}{m m m}{%
  \begin{exposition}%
  This chapter moves the development from #1 to #2. To get there, it studies: #3.%
  \end{exposition}%
}
